% CM splitting: the injection Cl(O_{K^+}) -> Cl(O_K) and h = h^+ h^-.

In this chapter $K$ denotes the $p$th cyclotomic field for an odd prime
$p$.  The goal is the factorisation $h = h^+ h^-$ of the class number: we
prove that the natural map on class groups induced by the inclusion
$\mathcal O_{K^+} \subseteq \mathcal O_K$ is injective, so that $h^+ \mid
h$ and the relative class number $h^- = h/h^+$ makes sense as a natural
number.  The key step is a descent statement for principal ideals
(\cref{thm:is-principal-after-extension}): an ideal of $\mathcal O_{K^+}$
that becomes principal in $\mathcal O_K$ was principal to begin with.

\section{The real subfield and the class numbers}

\begin{definition}\label{def:cm-real-subfield}
  \lean{NumberField.maximalRealSubfield}\leanok
For a number field $K$, let $K^+$ denote its maximal real subfield, the
largest subfield of $K$ all of whose complex embeddings are real.  When $K$
is a CM field, $K^+$ is the fixed field of complex conjugation, and
$[K : K^+] = 2$.  For $K = \Qzetap$ this is the usual maximal real
subfield $\Q(\zeta_p + \zeta_p^{-1})$.
\end{definition}

\begin{definition}\label{def:class-number-h}
  \lean{KummerCriterion.h}\leanok
For a number field $K$, the class number $h(K)$ is the cardinality of the
ideal class group of its ring of integers:
\[
  h(K)=\#\Cl{\mathcal O_K}.
\]
\end{definition}

\begin{definition}\label{def:class-number-hplus}
  \lean{KummerCriterion.hPlus}\leanok
  \uses{def:cm-real-subfield, def:class-number-h}
For a CM field $K$, the plus class number is the class number of the
maximal real subfield:
\[
  h^+(K)=h(K^+).
\]
\end{definition}

\section{Descent of principal ideals}

\begin{theorem}\label{thm:ring-of-integers-faithfully-flat}
  \lean{KummerCriterion.ringOfIntegers_faithfullyFlat_maximalRealSubfield}\leanok
  \uses{def:cm-real-subfield}
For every number field $L$, the ring $\mathcal O_L$ is faithfully flat over
$\mathcal O_{L^+}$.
\end{theorem}
\begin{proof}\leanok
$\mathcal O_L$ is a finitely generated torsion-free module over the
Dedekind domain $\mathcal O_{L^+}$, hence projective, hence flat.  For
faithfulness we use the criterion in terms of prime spectra: it suffices
that $\operatorname{Spec}\mathcal O_L \to \operatorname{Spec}\mathcal
O_{L^+}$ be surjective.  Given a prime $\mathfrak p$ of $\mathcal O_{L^+}$,
the extension $\mathcal O_{L^+} \subseteq \mathcal O_L$ is integral, so
lying-over produces a prime of $\mathcal O_L$ contracting to $\mathfrak p$.
\end{proof}

\begin{theorem}\label{thm:map-comap-ring-of-integers}
  \lean{KummerCriterion.map_comap_eq_ringOfIntegers}\leanok
  \uses{thm:ring-of-integers-faithfully-flat}
If $J$ is an ideal of $\mathcal O_{L^+}$, then extension to $\mathcal O_L$
followed by contraction is the identity:
\[
  (J\mathcal O_L)\cap\mathcal O_{L^+}=J.
\]
\end{theorem}
\begin{proof}\leanok
For a faithfully flat ring map $A \to B$, every ideal $J \subseteq A$
satisfies $JB \cap A = J$: the inclusion $J \subseteq JB \cap A$ is clear,
and the quotient map $A/J \to B/JB$ is injective by faithful flatness of
$B/JB$ over $A/J$.  Apply this with
\cref{thm:ring-of-integers-faithfully-flat}.
\end{proof}

\begin{theorem}\label{thm:is-principal-map-span-singleton}
  \lean{KummerCriterion.isPrincipal_of_map_eq_span_singleton_of_mem}\leanok
  \uses{thm:map-comap-ring-of-integers}
Let $I$ be an ideal of $\mathcal O_{K^+}$.  If $I\mathcal O_K$ is generated
by the image of an element $b_0\in\mathcal O_{K^+}$, then $I$ is principal,
generated by $b_0$.
\end{theorem}
\begin{proof}\leanok
Contract the equality $I\mathcal O_K=(b_0)$ back to $\mathcal O_{K^+}$.
The left-hand side contracts to $I$ by
\cref{thm:map-comap-ring-of-integers}.  For the right-hand side, the
contraction of $b_0\mathcal O_K$ contains $b_0\mathcal O_{K^+}$, and
conversely any element of $b_0 \mathcal O_K \cap \mathcal O_{K^+}$ is of
the form $b_0 c$ with $c \in K \cap \mathcal O_K = \mathcal O_{K^+}$
(using $b_0 \ne 0$; if $b_0 = 0$ then $I = 0$ is principal anyway).  Hence
$I=(b_0)$.
\end{proof}

\begin{theorem}\label{thm:ideal-map-conj}
  \lean{KummerCriterion.ideal_map_conj_eq}\leanok
  \uses{def:cm-real-subfield}
If $I$ is an ideal of $\mathcal O_{K^+}$, then the extended ideal
$I\mathcal O_K$ is fixed by complex conjugation.
\end{theorem}
\begin{proof}\leanok
The extension $I\mathcal O_K$ is generated by elements coming from $K^+$,
and those elements are fixed by conjugation.  Conjugation is a ring
automorphism of $\mathcal O_K$, so it maps $I\mathcal O_K$ onto the ideal
generated by the conjugates of the generators, which is $I\mathcal O_K$
itself.
\end{proof}

\begin{theorem}\label{thm:conj-generator-associated}
  \lean{KummerCriterion.conj_generator_associated}\leanok
  \uses{thm:ideal-map-conj}
Suppose $I \ne 0$ and $I\mathcal O_K=(a)$.  Then, by
\cref{thm:ideal-map-conj}, also $I\mathcal O_K=(\overline a)$, and the two
generators are associated: there is a unit $u\in\mathcal O_K^\times$ with
\[
  \overline a = u a.
\]
\end{theorem}
\begin{proof}\leanok
Two generators of the same nonzero principal ideal in a domain differ by a
unit: from $(a) = (\overline a)$ we get $\overline a = ua$ and $a = v
\overline a$ for some $u, v \in \mathcal O_K$, whence $a = vua$ and, since
$a \ne 0$, $vu = 1$.
\end{proof}

\begin{theorem}\label{thm:conj-unit-mul-one}
  \lean{KummerCriterion.conj_unit_mul_eq_one}\leanok
  \uses{thm:conj-generator-associated}
The unit $u$ obtained from $\overline a=ua$ in
\cref{thm:conj-generator-associated} satisfies
\[
  u\overline u=1.
\]
\end{theorem}
\begin{proof}\leanok
Apply complex conjugation to $\overline a=ua$: since conjugation is an
involution, $a = \overline u\,\overline a = \overline u u a$.  As $a\ne0$,
cancellation in the domain $\mathcal O_K$ gives $u\overline u=1$.
\end{proof}

\begin{theorem}\label{thm:antisymmetric-unit-classification}
  \lean{KummerCriterion.antisymmetric_unit_eq_neg_one_pow_mul_zeta_pow}\leanok
  \uses{thm:conj-unit-mul-one}
Let $K$ be the $p$th cyclotomic field with $p$ odd.  If
$u\in\mathcal O_K^\times$ satisfies $u\overline u=1$, then
\[
  u=(-1)^n\zeta_p^m
\]
for some integers $n,m$.
\end{theorem}
\begin{proof}\leanok
At every complex embedding $\sigma$ of $K$, conjugation corresponds to
complex conjugation of the image, so $|\sigma(u)|^2 = \sigma(u
\overline u) = 1$.  An algebraic integer all of whose conjugates have
absolute value $1$ is a root of unity (Kronecker's theorem), so $u$ is a
torsion unit.  The torsion subgroup of $\Qzetap^\times$ for $p$ odd is
generated by $-1$ and $\zeta_p$, which gives the displayed form.
\end{proof}

\begin{theorem}\label{thm:conj-fixed-associate}
  \lean{KummerCriterion.exists_conj_fixed_associate_of_classification}\leanok
  \uses{thm:antisymmetric-unit-classification}
In the situation of \cref{thm:conj-generator-associated}, the generator $a$
may be replaced by an associate $b$ such that
\[
  \overline b=b
  \qquad\text{and}\qquad
  (b)=(a).
\]
\end{theorem}
\begin{proof}\leanok
By \cref{thm:antisymmetric-unit-classification}, the unit $u =
\overline a / a$ has the form $(-1)^n \zeta_p^m$.  Since $p$ is odd, every
power of $\zeta_p$ is a square in $\langle\zeta_p\rangle$, and the sign can
be absorbed: write $u = w/\overline w$ for a suitable root of unity $w$
(concretely, a power of $\zeta_p$ chosen so that $w^2$ accounts for
$\zeta_p^m$, possibly times $-1$).  Then $b := w a$ satisfies
$\overline b = \overline w\,\overline a = \overline w u a = w a = b$,
and $b$ is an associate of $a$, so it generates the same ideal.
\end{proof}

\begin{theorem}\label{thm:mem-ring-of-integers-conj-self}
  \lean{KummerCriterion.mem_ringOfIntegers_of_conj_eq_self}\leanok
  \uses{def:cm-real-subfield}
If $b\in\mathcal O_K$ is fixed by complex conjugation, then it is the image
of an element of $\mathcal O_{K^+}$.
\end{theorem}
\begin{proof}\leanok
For a CM field, $K^+$ is exactly the fixed field of conjugation, so the
element $b$ lies in $K^+$.  Being integral over $\Z$, it lies in
$K^+ \cap \mathcal O_K = \mathcal O_{K^+}$.
\end{proof}

\begin{theorem}\label{thm:is-principal-after-extension}
  \lean{KummerCriterion.isPrincipal_of_isPrincipal_map_Kplus}\leanok
  \uses{thm:is-principal-map-span-singleton, thm:ideal-map-conj,
  thm:conj-generator-associated, thm:conj-unit-mul-one,
  thm:antisymmetric-unit-classification, thm:conj-fixed-associate,
  thm:mem-ring-of-integers-conj-self}
Let $K$ be the $p$th cyclotomic field, with $p$ odd.  If an ideal of
$\mathcal O_{K^+}$ becomes principal after extension to $\mathcal O_K$,
then it was already principal.
\end{theorem}
\begin{proof}\leanok
Let $I$ be such an ideal; we may assume $I \ne 0$.  Choose a generator $a$
of $I\mathcal O_K$.  The extended ideal is conjugation-stable
(\cref{thm:ideal-map-conj}), so \cref{thm:conj-generator-associated}
produces a unit $u$ with $\overline a = ua$, and
\cref{thm:conj-unit-mul-one} shows $u\overline u = 1$.  The classification
of such units (\cref{thm:antisymmetric-unit-classification}) lets us
replace $a$ by a conjugation-fixed associate $b$ with $(b) = I\mathcal O_K$
(\cref{thm:conj-fixed-associate}).  By
\cref{thm:mem-ring-of-integers-conj-self}, $b$ descends to an element $b_0
\in \mathcal O_{K^+}$, and \cref{thm:is-principal-map-span-singleton}
contracts the equality $I\mathcal O_K = (b_0)$ to $I = (b_0)$.
\end{proof}

\section{The relative class number}

\begin{definition}\label{def:class-group-map}
  \lean{KummerCriterion.classGroupMap}\leanok
  \uses{def:cm-real-subfield}
The inclusion $\mathcal O_{K^+}\subseteq\mathcal O_K$ induces, by extension
of fractional ideals, the class-group homomorphism
\[
  \Cl{\mathcal O_{K^+}}\longrightarrow \Cl{\mathcal O_K}.
\]
\end{definition}

\begin{theorem}\label{thm:class-group-map-injective}
  \lean{KummerCriterion.classGroupMap_injective}\leanok
  \uses{def:class-group-map, thm:is-principal-after-extension}
For an odd prime cyclotomic field, the map
\[
  \Cl{\mathcal O_{K^+}}\longrightarrow \Cl{\mathcal O_K}
\]
is injective.
\end{theorem}
\begin{proof}\leanok
A group homomorphism is injective when its kernel is trivial.  If an ideal
class of $\mathcal O_{K^+}$ maps to the trivial class, a representative
ideal becomes principal after extension to $\mathcal O_K$; by
\cref{thm:is-principal-after-extension} the representative was already
principal, so the class is trivial.
\end{proof}

\begin{theorem}\label{thm:hplus-divides-h}
  \lean{KummerCriterion.hPlus_dvd_h}\leanok
  \uses{def:class-number-h, def:class-number-hplus,
  thm:class-group-map-injective}
The plus class number divides the full class number:
\[
  h^+(K)\mid h(K).
\]
\end{theorem}
\begin{proof}\leanok
By \cref{thm:class-group-map-injective}, $\Cl{\mathcal O_{K^+}}$ embeds as
a subgroup of the finite group $\Cl{\mathcal O_K}$, and Lagrange's theorem
gives the divisibility of cardinalities.
\end{proof}

\begin{definition}\label{def:class-number-hminus}
  \lean{KummerCriterion.hMinus}\leanok
  \uses{thm:hplus-divides-h}
The relative class number is the natural-number quotient
\[
  h^-(K)=h(K)/h^+(K),
\]
which is well defined by \cref{thm:hplus-divides-h}.
\end{definition}

\begin{theorem}\label{thm:h-factorisation}
  \lean{KummerCriterion.h_eq_hPlus_mul_hMinus}\leanok
  \uses{def:class-number-hminus, thm:hplus-divides-h}
The class number factors as
\[
  h(K)=h^+(K)\,h^-(K).
\]
\end{theorem}
\begin{proof}\leanok
Multiply the defining quotient of \cref{def:class-number-hminus} by
$h^+(K)$; this is legitimate in $\N$ because $h^+(K) \mid h(K)$
(\cref{thm:hplus-divides-h}) and $h^+(K) \ne 0$.
\end{proof}

\begin{theorem}\label{thm:dvd-h-iff-dvd-hminus-from-plus-to-minus}
  \lean{KummerCriterion.dvd_h_iff_dvd_hMinus_of_dvd_hPlus_imp}\leanok
  \uses{thm:h-factorisation}
Assume that a separate argument proves the implication
\[
  p\mid h^+(K)\Longrightarrow p\mid h^-(K).
\]
Then
\[
  p\mid h(K)\Longleftrightarrow p\mid h^-(K).
\]
\end{theorem}
\begin{proof}\leanok
Write $h(K) = h^+(K)\,h^-(K)$ by \cref{thm:h-factorisation}.  If $p \mid
h(K)$, then by primality $p\mid h^+(K)$ or $p\mid h^-(K)$, and the first
alternative implies the second by hypothesis; either way $p \mid h^-(K)$.
Conversely $h^-(K)$ divides $h(K)$, so $p \mid h^-(K)$ implies
$p \mid h(K)$.
\end{proof}
